\section*{Consignes}
\begin{itemize}
  \item Durée totale : \textbf{1h30}.
  \item Un squelette \texttt{.c} est fourni par exercice : complétez-le, ne renommez pas les fichiers.
  \item Chaque programme doit compiler sans erreur et sans avertissement.
  \item Attention, le résultat doit être affiché exactement comme indiqué dans les exemples, sinon la correction automatique ne pourra pas le valider. Les lignes surlignées ci-dessous sont celles effectivement vérifiées par check50 : le reste (invites de saisie) n'est donné qu'à titre indicatif.
  \item En cas d'erreur, c'est à vous de vous adapter au check50, non l'inverse. Vous pouvez relancer le check50 autant de fois que nécessaire.
  \item A la fin de l'épreuve, faites un submit50 de votre dossier :
  \begin{lstlisting}[language=bash,style=console]
submit50 IUT-GEII-Annecy/exercices/2026/info1/eval_pratique_1
  \end{lstlisting}
\end{itemize}

\begin{UPSTIexercice}{Variables et affichage (10 min)}
  \textbf{Fichier :} \texttt{variables.c}
  \begin{itemize}
    \item Demander à l'utilisateur une intensité et une tension.
    \item Calculer la puissance électrique correspondante.
    \item Afficher le résultat en W, avec deux chiffres après la virgule.
  \end{itemize}

  \textbf{Exemple :}
  \begin{lstlisting}[language=bash,style=console,escapeinside={(*}{*)}]
Intensité (A) : 4
Tension (V) : 230
(*\colorbox{yellow!40}{\ttfamily Puissance : 920.00 W}*)
  \end{lstlisting}
\end{UPSTIexercice}

\begin{UPSTIexercice}{Condition simple (10 min)}
  \textbf{Fichier :} \texttt{condition.c}
  \begin{itemize}
    \item Demander une tension à l'utilisateur.
    \item Le programme devra afficher \texttt{Tension nominale} entre 220 V et 230 V inclus. Au dessus, afficher \texttt{Surtension}, en dessous, afficher \texttt{Sous-tension}.
  \end{itemize}

  \textbf{Exemple :}
  \begin{lstlisting}[language=bash,style=console,escapeinside={(*}{*)}]
Tension (V) : 235
(*\colorbox{yellow!40}{\ttfamily Surtension}*)
  \end{lstlisting}
\end{UPSTIexercice}

\begin{UPSTIexercice}{Compte à rebours (15 min)}
  \textbf{Fichier :} \texttt{compte\_a\_rebours.c}
  \begin{itemize}
    \item Demander un nombre de départ à l'utilisateur.
    \item Afficher le compte à rebours, une valeur par ligne, jusqu'à 1.
    \item Afficher ensuite \texttt{"Décollage !"}.
    \item Aucune temporisation n'est demandée.
  \end{itemize}

  \textbf{Exemple :}
  \begin{lstlisting}[language=bash,style=console,escapeinside={(*}{*)}]
n : 5
(*\colorbox{yellow!40}{\ttfamily 5}*)
(*\colorbox{yellow!40}{\ttfamily 4}*)
(*\colorbox{yellow!40}{\ttfamily 3}*)
(*\colorbox{yellow!40}{\ttfamily 2}*)
(*\colorbox{yellow!40}{\ttfamily 1}*)
(*\colorbox{yellow!40}{\ttfamily Décollage !}*)
  \end{lstlisting}
\end{UPSTIexercice}

\begin{UPSTIexercice}{Boucle avec accumulateur (15 min)}
  \textbf{Fichier :} \texttt{accumulateur.c}
  \begin{itemize}
    \item Demander à l'utilisateur un nombre de mesures.
    \item Saisir ensuite ces mesures de tension, une à une.
    \item Calculer et afficher leur moyenne.
    \begin{itemize}
      \item Afficher le résultat en V, avec deux chiffres après la virgule. Ce résultat sera précédé de la phrase \texttt{"Moyenne : "} et suivi d'un retour à la ligne.
    \end{itemize}
  \end{itemize}

  \textbf{Exemple :}
  \begin{lstlisting}[language=bash,style=console,escapeinside={(*}{*)}]
Nombre de mesures : 3
Tension 1 (V) : 220
Tension 2 (V) : 225
Tension 3 (V) : 230
(*\colorbox{yellow!40}{\ttfamily Moyenne : 225.00 V}*)
  \end{lstlisting}
\end{UPSTIexercice}

\begin{UPSTIexercice}{Fonction avec paramètres et retour (15 min)}
  \textbf{Fichier :} \texttt{diviseur\_tension.c}
  \begin{itemize}
    \item Demander à l'utilisateur une tension d'entrée $v_{in}$ et deux résistances $r_1$, $r_2$ d'un pont diviseur de tension.
    \item Calculer la tension de sortie $v_{out} = v_{in} \times r_2 / (r_1 + r_2)$ à l'aide d'une fonction dédiée.
    \item Afficher le résultat en V, avec deux chiffres après la virgule. Ce résultat sera précédé de la phrase \texttt{"Tension de sortie : "} et suivi d'un retour à la ligne.
  \end{itemize}

  \textbf{Exemple :}
  \begin{lstlisting}[language=bash,style=console,escapeinside={(*}{*)}]
Vin (V) : 12
R1 (ohm) : 1000
R2 (ohm) : 2000
(*\colorbox{yellow!40}{\ttfamily Tension de sortie : 8.00 V}*)
  \end{lstlisting}
\end{UPSTIexercice}

\begin{UPSTIexercice}{Synthèse : boucle + condition + fonction (20 min)}
  \textbf{Fichier :} \texttt{synthese.c}
  \begin{itemize}
    \item Demander à l'utilisateur une tolérance (une valeur minimale et une valeur maximale), puis un nombre de mesures.
    \item Saisir ces mesures une à une et compter, à l'aide d'une fonction dédiée, celles qui sont strictement hors de la tolérance.
    \item Afficher le nombre total de mesures hors tolérance. Ce résultat (un entier) sera précédé de la phrase \texttt{"Mesures hors tolérance : "} et suivi d'un retour à la ligne.
  \end{itemize}

  \textbf{Exemple :}
  \begin{lstlisting}[language=bash,style=console,escapeinside={(*}{*)}]
Min : 220
Max : 230
Nombre de mesures : 4
Mesure 1 : 215
Mesure 2 : 225
Mesure 3 : 235
Mesure 4 : 228
(*\colorbox{yellow!40}{\ttfamily Mesures hors tolérance : 2}*)
  \end{lstlisting}
\end{UPSTIexercice}

\begin{UPSTIexercice}{Bonus -- Récursivité : PGCD (optionnel, hors barème et hors temps de base)}
  \textbf{Fichier :} \texttt{pgcd.c}

  \textbf{Ne compte pas dans la note ni dans les 85 minutes de base.} À traiter uniquement si le reste est terminé.
  \begin{itemize}
    \item Demander à l'utilisateur deux entiers.
    \item Calculer leur PGCD à l'aide d'une fonction \textbf{récursive} (algorithme d'Euclide).
    \item Afficher le résultat (un entier). Ce résultat sera précédé de la phrase \texttt{"PGCD : "} et suivi d'un retour à la ligne.
  \end{itemize}

  \textbf{Exemple :}
  \begin{lstlisting}[language=bash,style=console,escapeinside={(*}{*)}]
a : 48
b : 18
(*\colorbox{yellow!40}{\ttfamily PGCD : 6}*)
  \end{lstlisting}
\end{UPSTIexercice}
